\documentclass[11pt]{article}
\usepackage{amsmath,hyperref}
\title{Produce the scope at the gate that owns later changes}
\date{}
\begin{document}\maketitle

\paragraph{Result and scope.}
A stored statement that an older operation leaves $Y$ alone need not
remain true when a later operation acquires $Y$. In the authored warehouse
simulator, an ordered source gate extracts currently pending resources
from its actual command queue and checks the live claim against later queue
changes. The earlier producer is no longer a presumed pass bit: this
fixture produces local scope records from qualified native executions
and enforces the modeled amendment law. It is not a claim about arbitrary
controllers or deployed warehouse sensing.

\paragraph{What the source checks.}
Reuse the exact old deliveries and live buffer prefix of
\href{https://github.com/eternal-november/eos-research-programmes/pull/586}
{the predecessor-scope case}. The source's accepted suffix consists of
typed native operations with an actor, intended movement and before and
after cell and carried-item values. For each old and new step used in
the experiment the fixture checks native actor position, heading and
carried item; before pickup it checks the expected physical shelf at
the site. This is an authored native-state observation interface, not
hardware qualification. Stationary footprints on the routes in this
fixed fixture and the prescribed access law are held fixed.

For each pending old operation the extracted resources include the
actor, both cells and both applicable carried-item identities. Their
union, revision and consumed step count form a read record. Every
accepted old queue installation or edit must pass through the same
ordered source gate. Successful acquisition rechecks the current
record, then holds the new job's footprint until the actual terminal
effect is checked. Ignoring other authorities, hidden queue writers
or other work that can reach these same cells would invalidate the
coverage argument. There is no cryptographic signature or network
round-trip model. The full pending record stays inside the local
fixture adapter. There is no claim of privacy if it is sent elsewhere.

\paragraph{Distinct native cuts.}
\begin{itemize}
\item A conflicting edit that wins \emph{before} acquisition makes
the earlier source record stale. The new read remains conflicting
while the old exchange still needs the same shelf. When the old
partner physically clears, the new move can instead follow the
revised eligible bay.
\item The same edit offered \emph{after} acquisition is refused.
A disjoint parking turn is accepted in another run while the
claim stays live, so this is not a freeze on all old changes.
\item During the actual exchange, at tick 8 the older carrier
has unloaded $Y$. No remaining operation touches the \emph{asset}
$Y$, but the older carrier still occupies its cell and must clear
it. Actor and path scope therefore prohibit acquisition at that
cut. At tick 12 the carrier really is clear and the source-derived
suffix is disjoint. Actual new relocation ends at tick 18, before
the whole older exchange finishes at tick 28.
\item Delaying the tick-12 record by one additional old step
requires rechecking. The original revision's queue position is
no longer current even if the resource set remains disjoint.
The same source supplies a fresh read at tick 13 and the
physical move ends at 19. The direct handoff could also refresh
immediately when authorized; no round-trip speed advantage is claimed.
\end{itemize}

All tested combinations preserve the old effects, the licensed
new relocation, its subsequent native goal receipt and independent
use. Actual release checks that the new holder has unloaded at the
qualified refuge and cleared storage of the old footprint. The
subsequent goal in this producer demonstration follows the bay
actually granted by the fixture; it does not replace the fixed
customer-path test in the predecessor experiment.

\paragraph{Proof and comparison.}
The local invariant is that the currently held footprint is disjoint
from the accepted old suffix. Deleting a verified head preserves
disjointness. Installing an amended suffix must recheck against a
held footprint; retaining the current one on refusal preserves
the invariant. The three Lean statements check these set arguments,
not the native abstraction or a distributed serialization theorem.
A closure covering only assets cannot substitute for the tested
spatial cuts. This is a source-local use of the idea of witnessed
phase release, not a claim of advantage over a matched
\href{https://cs.nyu.edu/~shasha/papers/chopping.pdf}
{transaction-chopping} implementation.

The comparison uses direct same-gate lookup of each command's
fields and the same actuation and observation. It agrees.
Ten native cases verify 234 old command steps, scan 1,042 old
command entries across scope checks and scan 146 proposed
entries in edit checks; these are operation counts, not execution
time or throughput. Single owner, one active new claim and a closed
old-opcode set are explicit. The result is not an independent
architecture transfer, fault disposition or acceptance criterion.
\end{document}
